\SourceLeaf{323}
\chapter*{ZUR THEORIE DER BIQUADRATISCHEN RESTE.}
\pdfbookmark[0]{Zur Theorie der biquadratischen Reste}{nachlass-biquadratic}
\renewcommand{\BookRunningTitle}{Zur Theorie der biquadratischen Reste}
\BookRule{.18\textwidth}
\begin{center}\textbf{[I.]}\end{center}
\pdfbookmark[1]{I}{nachlass-biquadratic-I}
\renewcommand{\Article}[1]{\par\medskip\Needspace{4\baselineskip}\pdfbookmark[2]{Art. #1}{nachlass-biquadratic-I-art-#1}\centerline{\textbf{#1.}}\nobreak\smallskip}
\Article{1}
Wir erweitern das Gebiet der höheren Arithmetik, indem wir darin auch die komplexen Größen aufnehmen. In dieser Untersuchung nennen wir jede Größe $x+iy$ eine komplexe ganze Zahl, wenn $x,y$ reelle ganze Zahlen sind.

\Article{2}
Die unendlich vielen komplexen ganzen Zahlen lassen sich am bequemsten durch Punkte in einer unbegrenzten Ebene anschaulich darstellen; wir nennen den Punkt mit der Abszisse $x$ und der Ordinate $y$ den Punkt $x+iy$; alle Punkte, die ganze Zahlen darstellen, sollen ganzzahlige Punkte heißen.

\Article{3}
Um eine bestimmte Festlegung zu treffen, sollen die Abszissen stets nach links positiv und die Ordinaten nach oben positiv sein.

\Article{4}
Die vom Punkt $x+iy$ zum Punkt $x'+iy'$ gezogene gerade Linie soll die gerade Linie $(x+iy,x'+iy')$ heißen; dabei berücksichtigen wir, wo es darauf ankommt, zugleich die Richtung und unterscheiden daher die gerade Linie $x+iy,x'+iy'$ von $x'+iy',x+iy$.

\SourceLeaf{324}
\Article{5}
Der Kürze halber wollen wir komplexe Größen wie $x+iy$ auch durch einen einzigen Buchstaben bezeichnen, etwa $z$.

\Article{6}
Die von den geraden Linien $zz',z'z'',z''z'''\ldots z^{n-1}z^n,z^nz$ begrenzte Figur nennen wir die Figur $zz'z''z'''\dotstwo z^n$. Wir schließen den Fall nicht aus, dass einige dieser Linien einander schneiden.

\Article{7}
Mit $S(z,z',z''\ldots z^n)$ bezeichnen wir allgemein die Summe von so vielen reellen ganzen Zahlen, wie ganzzahlige Punkte innerhalb der Figur liegen. Dabei setzen wir für jeden Punkt, den die Grenzlinie der Figur einmal, zweimal, dreimal usw. umläuft, die Zahl $\pm1$, $\pm2$, $\pm3$ usw.; die oberen Zeichen gelten, wenn die Grenzlinie den Punkt so umgibt, dass er auf der rechten Seite der Figur liegt, die unteren im entgegengesetzten Fall. Schneiden sich also keine Seiten der Figur, so ist $S(z,z',z''\dotstwo)$ die Anzahl der Punkte innerhalb der Figur, positiv oder negativ genommen.

\Article{8}
Offenbar gilt stets
\[
\begin{aligned}
S(z,z',z''\ldots z^n)&=S(z',z'',z'''\dotstwo z^n,z)\\
&=S(z'',z'''\dotstwo z')\text{ usw.}\\
&=-S(z^n,z^{n-1}\dotstwo z'',z',z)\\
&=-S(z^{n-1}\pdot z^{n-2}\dotstwo z',z,z^n)\text{ usw.}
\end{aligned}
\]

\Article{9}
Wie die auf der Grenzlinie selbst liegenden Punkte zu behandeln sind, muss noch genauer festgelegt werden. In vielen Fällen können auf der Grenzlinie überhaupt keine ganzzahligen Punkte liegen; dann ist keine Festlegung nötig. Liegen aber auf der Grenzlinie $zz'$ solche Punkte, so zeigen wir durch ein zwischen $z$ und $z'$ eingefügtes $+$ an, dass diese Punkte so zu behandeln sind, als lägen sie rechts von der Grenzlinie, und durch ein $-$, als lägen sie links. Wir werden auch gelegentlich ein $0$ oder $\frac12$ einfügen, um anzuzeigen, dass sie auf
\SourceLeaf{325}
jeder Seite gar nicht oder nur mit ihrem halben Wert zu berücksichtigen sind. Ist einer der Punkte $z,z',z''$ usw. selbst ein ganzzahliger Punkt, so wird er, sofern nicht ausdrücklich das Gegenteil gesagt wird, überhaupt nicht mitgezählt, außer soweit er zugleich auch als Punkt, der kein Eckpunkt ist, in Betracht kommt.

\Article{10}
Sätze. Werden alle $z,z',z''$ usw. um dieselbe ganze Zahl vermehrt, so bleibt $S$ unverändert.

Wird $i$ durch $-i$ und jedes Verbindungszeichen durch sein entgegengesetztes ersetzt, so ändert $S$ nur sein Vorzeichen.
\[
\begin{aligned}
S(z,z',z''\ldots z^n)={}&S(z,u,u'\dotstwo u^n,z',z'',z^n)\\
&-S(z,u,u'\dotstwo u^n,z')\\
={}&S(z,u,u',u''\dotstwo u^n,z^m,z^{m+1}\dotstwo z^n)\\
&-S(z,u,u',u''\dotstwo u^n,z^m,z^{m-1}\dotstwo z',z)
\end{aligned}
\]
Dabei müssen die Verbindungszeichen einander entsprechen, zwischen den rückwärts laufenden Gliedern jedoch umgekehrt werden.

Ist $\zeta$ eine ganze Zahl $=a+bi$, so gilt, wenn die gegenüberliegenden Verbindungszeichen entgegengesetzt sind,
\[
S(z,z',z'+\zeta,z+\zeta)=[bx'-ay']-[bx-ay]
\]
Ist hierbei $bx'-ay'$ selbst eine ganze Zahl, so ist diese für $[bx'-ay']$ zu nehmen, wenn das Verbindungszeichen zwischen $z'$ und $z'+\zeta$ $+$ ist, dagegen $1$ oder $\frac12$ weniger, wenn dieses Zeichen $-$ oder $\frac12$ ist; für $bx-ay$ gilt das Umgekehrte.

Die Formel gilt nur, wenn $a$ und $b$ keinen gemeinsamen Teiler haben; ist ihr größter gemeinsamer Teiler $=h$, so ist stattdessen
\[
h\left[\frac{bx'-ay'}h\right]-h\left[\frac{bx-ay}h\right]
\]
zu nehmen.
